\documentclass[11pt,a4paper]{article}
\usepackage[utf8]{inputenc}
\usepackage[spanish]{babel}
\usepackage{amsmath,amsfonts,amssymb}
\usepackage{geometry}
\geometry{top=2.5cm, bottom=2.5cm, left=2.5cm, right=2.5cm}
\usepackage{hyperref}
\usepackage{booktabs}
\usepackage{tabularx}
\usepackage{xcolor}
\usepackage{tcolorbox}
\usepackage{enumitem}

% Definición de colores institucionales
\definecolor{primary}{RGB}{15, 32, 67}
\definecolor{secondary}{RGB}{0, 114, 206}
\definecolor{lightgray}{RGB}{245, 247, 250}

\hypersetup{
    colorlinks=true,
    linkcolor=secondary,
    urlcolor=secondary,
    pdftitle={Portafolio de Evidencias No. 2 - StatCredit AI}
}

\begin{document}

% ---------------------------------------------------------
% PORTADA INSTITUCIONAL
% ---------------------------------------------------------
\begin{titlepage}
    \centering
    \vspace*{1cm}
    {\Huge \textbf{UNIVERSIDAD EXTERNADO DE COLOMBIA}}\\[0.4cm]
    {\Large \textbf{PREGRADO EN CIENCIA DE DATOS --- EMPRENDIMIENTO}}\\[1.8cm]
    
    \rule{\linewidth}{0.8mm}\\[0.4cm]
    {\Huge \textbf{PORTAFOLIO DE EVIDENCIAS NO. 2}}\\[0.3cm]
    {\Large \textbf{Consolidado de Entregables: Tareas 1, 2 y 3}}\\[0.4cm]
    \rule{\linewidth}{0.8mm}\\[1.8cm]
    
    \begin{tcolorbox}[colback=primary!5,colframe=primary]
        \centering \Huge
        \textbf{StatCredit AI}
    \end{tcolorbox}
    
    \vfill
    
    \begin{minipage}{0.8\textwidth}
        \centering
        \textbf{\large Integrantes (Grupo 5):}\\[0.3cm]
        Santiago Sandoval \quad\vert\quad Sebastián Ramos\\[0.15cm]
        Julián Duarte \quad\vert\quad Tomás Rincón \quad\vert\quad Julián Jiménez
    \end{minipage}
    
    \vfill
    {\large \textbf{18 de Septiembre de 2026 --- Bogotá D.C., Colombia}}
\end{titlepage}

\newpage
\tableofcontents
\newpage

% ---------------------------------------------------------
% SECCIÓN 1: TAREA 1
% ---------------------------------------------------------
% ---------------------------------------------------------
% SECCIÓN 1: TAREA 1
% ---------------------------------------------------------
\section{Tarea 1: Taller de Ideación y Validación con Expertos}

\begin{tcolorbox}[colback=primary!5,colframe=primary,title=\textbf{¿De qué trata este entregable?}]
En esta primera tarea validamos en campo la hipótesis del problema de \textbf{StatCredit AI} mediante entrevistas cualitativas en profundidad con expertos del sector financiero en Colombia (un ejecutivo de riesgo crediticio de consumo y una experta en regulación actuarial SARC). El objetivo es corroborar si el punto ciego de los independientes y la necesidad de explicabilidad en Inteligencia Artificial representan dolores reales y vigentes en la banca.
\end{tcolorbox}

\subsection{El Problema de Forma Sencilla y Reencuadre}

Muchos trabajadores independientes solventes en Colombia (comerciantes urbanos, profesionales autónomos y micronegociantes) poseen una intensa actividad financiera diaria en billeteras digitales (Nequi, Daviplata) y facturación electrónica. Sin embargo, al solicitar crédito formal, son declinados por carecer de historial crediticio bancario o certificaciones tradicionales de nómina.

\textbf{Nuestro reencuadre del problema (Bianchi \& Verganti, 2021):}
El problema no es que los independientes carezcan de capacidad de repago. El verdadero problema es que los procesos tradicionales de evaluación no cuentan con la infraestructura analítica para ingestar e interpretar esa información transaccional observable. En StatCredit AI no buscamos cambiar a las personas, sino cambiar la forma en que se mide quién puede pagar un crédito.

\subsection{Caracterización de los Expertos Entrevistados}

Para validar estas hipótesis, realizamos entrevistas cualitativas estructuradas a dos perfiles estratégicos:

\begin{enumerate}[leftmargin=*]
    \item \textbf{Diego Ramos (Experto en Riesgo Crediticio):} Ejecutivo de crédito de consumo con años de experiencia directa en la evaluación operacional y comités de aprobación crediticia.
    \item \textbf{Paola Cárdenas (Experta en Regulación SARC):} Analista de crédito microempresarial y especialista en cumplimiento normativo de la Superintendencia Financiera de Colombia (SARC).
\end{enumerate}

\subsection{Transcripción y Resultados de las Entrevistas de Campo}

\subsubsection{Entrevista 1: Diego Ramos (Ejecutivo de Crédito de Consumo)}

\begin{tcolorbox}[colback=lightgray,colframe=primary,title=\textbf{Pregunta 1: Rechazo de independientes por historial en buró}]
\textit{¿Ustedes rechazan mucha gente independiente por lo del tema de DataCrédito y esas cosas?}
\tcblower
\textbf{Respuesta de Diego Ramos:}  
«Sí, muchísimo, más de lo que se cree normalmente. Yo llevaba años en eso y apenas el sistema ve que el cliente no tiene un puntaje bueno en el buró, ni siquiera llega a que un analista lo mire con calma. Y muchas veces ese señor sí tiene con qué pagar por lo que puede decirnos más personalmente, solo que nunca tuvo una tarjeta de crédito o nunca sacó un préstamo formal, entonces para el sistema es como si no existiera.»
\end{tcolorbox}

\begin{tcolorbox}[colback=lightgray,colframe=primary,title=\textbf{Pregunta 2: Capacidad de procesamiento manual de extractos digitales}]
\textit{¿Y si la persona le muestra los movimientos de Nequi o las facturas, ustedes tienen cómo revisar eso rápido?}
\tcblower
\textbf{Respuesta de Diego Ramos:}  
«No, y ahí es donde hay un problema ya que eso tocaría revisarlo a mano, uno abriendo el PDF del extracto, contando transacción por transacción, sacando el promedio en Excel y así uno alcanza a sacar máximo 15 o 20 casos al día por persona. Pero no es que no queramos usar esa información, es que no damos abasto para procesarla cuando hay tantas personas que manejan el dinero de esa manera. Si a mí me hubieran dado algo automático que me diera bien la información como lo que me cuentas que quieren hacer, eso me cambiaría la manera de cómo se rechazaría a los clientes.»
\end{tcolorbox}

\begin{tcolorbox}[colback=lightgray,colframe=primary,title=\textbf{Pregunta 3: Comportamiento de pago de independientes vs. asalariados}]
\textit{¿Un independiente que mueve plata todos los días es más moroso que uno con sueldo fijo?}
\tcblower
\textbf{Respuesta de Diego Ramos:}  
«No, y digamos que eso lo sabemos todos los que trabajan en esto, aunque no se mencione mucho. Yo he visto casos de gente que tiene su negocio como yo y le entra plata todos los días, y hay más organizados que muchos asalariados que conozco. Lo que pasa es que a uno en el banco lo criaron con el cuento de que 'nómina fija es igual a seguridad', y es como que ese chip es difícil de cambiar para que la gente de un día para otro revise otras perspectivas.»
\end{tcolorbox}

\begin{tcolorbox}[colback=lightgray,colframe=primary,title=\textbf{Pregunta 4: Miedo a la aprobación sin desprendible de nómina}]
\textit{¿Y qué tanto miedo le tienen ustedes a aprobar un crédito sin el desprendible de nómina de toda la vida?}
\tcblower
\textbf{Respuesta de Diego Ramos:}  
«El miedo no es nuestro sino del banco, porque muchas personas nos cuentan sus historias y de manera hablada catalogan para poder entrar a revisión pero sin las herramientas deja de ser posible. Si yo apruebo algo raro y después algo sale mal, quien tiene que ir a dar la cara en el comité de auditoría soy yo. Entonces pues uno prefiere por si las moscas, decir que no de una vez, así uno sepa que se le está negando el préstamo a alguien que seguramente sí paga.»
\end{tcolorbox}

\subsubsection{Entrevista 2: Paola Cárdenas (Analista de Crédito Microempresarial \& Experta SARC)}

\begin{tcolorbox}[colback=lightgray,colframe=secondary,title=\textbf{Pregunta 5: Requerimientos de la Superintendencia Financiera (SARC)}]
\textit{¿Qué les exige la Superfinanciera cuando un banco quiere usar algoritmos para dar créditos?}
\tcblower
\textbf{Respuesta de Paola Cárdenas:}  
«Lo que a uno le exigen de fondo con el tema del SARC es que la entidad pueda explicar en cualquier momento por qué se tomó tal decisión de crédito. Entonces a uno no le sirve decir 'es que el modelo dijo que no' y ya. Da igual si es una fórmula de estadística o un algoritmo de inteligencia artificial, la responsabilidad de explicarlo siempre es de la entidad, nunca del programa ni de quien se lo vendió.»
\end{tcolorbox}

\begin{tcolorbox}[colback=lightgray,colframe=secondary,title=\textbf{Pregunta 6: Tranquilidad regulatoria con reportes de explicabilidad (XAI/SHAP)}]
\textit{¿Y si nuestra herramienta les entrega un reporte fácil de entender explicando por qué la IA decidió cada cosa, eso les da tranquilidad para una auditoría?}
\tcblower
\textbf{Respuesta de Paola Cárdenas:}  
«Sí, digamos que si tienen un pleno soporte que el banco acepte se podría hacer ya que nos autorizarían a poder utilizarlo como un complemento, pero sin ese reportecito que explique qué pesó a favor y qué en contra, ningún oficial de cumplimiento se va a poner a defender esa decisión frente a un auditor, así el modelo sea buenísimo o ustedes tengan mucha información. Así que con esa explicación clara, la cosa deja de sonar a caja negra sospechosa y uno sí la puede defender más tranquilo, o en caso tal no toda la responsabilidad cae sobre el que acepta el crédito.»
\end{tcolorbox}

\subsection{Síntesis de Hallazgos y Validación de Hipótesis}

Las entrevistas de campo aportan evidencias cualitativas directas que respaldan de forma contundente la propuesta de valor de \textbf{StatCredit AI}:

\begin{table}[h!]
\centering
\small
\begin{tabularx}{\textwidth}{p{3.5cm} X p{4.5cm}}
\toprule
\textbf{Hipótesis Evaluada} & \textbf{Evidencia Cualitativa Obtenida en Campo} & \textbf{Decisión de Diseño en StatCredit AI} \\
\midrule
\textbf{1. Fricción por Ingesta Manual} & Diego Ramos confirma que los analistas solo pueden procesar 15-20 extractos al día manualmente en Excel. & Desarrollo de la API REST para ingesta y cálculo automatizado en $<2$ segundos. \\
\midrule
\textbf{2. Falsos Negativos y Sesgo Cultural} & Diego evidencia que el sesgo de "nómina fija igual a seguridad" causa rechazo masivo de independientes solventes. & Implementación de modelos de Inferencia Causal para evaluar estabilidad en flujos estacionales. \\
\midrule
\textbf{3. Miedo al Comité y Explicabilidad SARC} & Paola Cárdenas señala que ningún oficial defenderá un modelo de "caja negra" ante la Superfinanciera (Circular 026). & Inclusión nativa de la matriz de explicabilidad SHAP en la respuesta JSON y el Dashboard. \\
\bottomrule
\end{tabularx}
\caption{Matriz de validación de hipótesis a partir de los testimonios de los expertos.}
\end{table}

\subsection{Conclusión de la Tarea 1}
La validación con los expertos Diego Ramos y Paola Cárdenas confirma que el dolor en el otorgamiento de crédito a independientes no se debe a desinterés institucional, sino a la falta de herramientas tecnológicas automatizadas de ingesta transaccional y reportes explicables adaptados a las exigencias del SARC. \textbf{StatCredit AI resuelve exactamente esta fricción operativa y regulatoria.}

\newpage
% ---------------------------------------------------------
% SECCIÓN 2: TAREA 2
% ---------------------------------------------------------
\section{Tarea 2: Experimentos para Medir y Probar la Idea}

\begin{tcolorbox}[colback=primary!5,colframe=primary,title=\textbf{¿De qué trata este entregable?}]
En esta segunda tarea proponemos \textbf{12 experimentos prácticos} organizados en 6 preguntas clave (2 experimentos por pregunta). Como estudiantes de Ciencia de Datos, el objetivo es aplicar herramientas de análisis de datos, pruebas de usuario y mediciones simples para comprobar si las ideas de \textbf{StatCredit AI} funcionan antes de gastar tiempo o dinero construyendo cosas que nadie vaya a usar.
\end{tcolorbox}

\subsection{Matriz de Experimentos}

\subsubsection{Pregunta 1: ¿Quién es tu cliente?}

\begin{tcolorbox}[colback=lightgray,colframe=secondary,title=\textbf{Experimento 1.1: Agrupar empresas para encontrar nuestro primer cliente}]
\begin{itemize}[leftmargin=*]
    \item \textbf{Lo que queremos saber:} ¿Cuáles entidades financieras en Colombia tienen más urgencia de dar crédito a trabajadores independientes?
    \item \textbf{Cómo lo probamos (Ciencia de Datos):} Tomamos datos públicos de entidades financieras en Colombia y aplicamos algoritmos de agrupación (\textit{clustering} como $k$-means). Evaluamos dos variables principales: el volumen de solicitudes de crédito rechazadas y su \textbf{flexibilidad tecnológica}. Entendemos por \textit{flexibilidad tecnológica} aquellas entidades (como fintechs o neobancos) que ya cuentan con plataformas digitales activas y arquitecturas basadas en APIs abiertas, a diferencia de los bancos tradicionales que aún dependen de procesos manuales en papel o sistemas legados cerrados.
    \item \textbf{Por qué importa la métrica:} Si encontramos un grupo claro de al menos 30 o 40 entidades con alta flexibilidad tecnológica y altos rechazos a independientes, sabremos exactamente a quiénes ofrecer primero nuestro servicio (\textit{Beachhead Market}).
\end{itemize}
\end{tcolorbox}

\begin{tcolorbox}[colback=lightgray,colframe=secondary,title=\textbf{Experimento 1.2: Prueba A/B en página web para ver qué mensaje llama más la atención}]
\begin{itemize}[leftmargin=*]
    \item \textbf{Lo que queremos saber:} ¿A los directores de riesgo les atrae más el mensaje de \textit{"reducir la mora en préstamos a independientes"} o el de \textit{"cumplir con las normas de Open Finance"}?
    \item \textbf{Cómo lo probamos:} Diseñamos una página web sencilla con dos versiones del mensaje principal (Prueba A/B) y medimos en cuál hace clic más gente interesada.
    \item \textbf{Por qué importa la métrica:} Si la versión de \textit{"reducir la mora"} logra el doble de clics y registros para demostraciones, confirmaremos que los tomadores de decisión priorizan el impacto financiero directo (no perder dinero) sobre el simple cumplimiento de normas regulatorias.
\end{itemize}
\end{tcolorbox}

\subsubsection{Pregunta 2: ¿Qué puedes hacer por tu cliente?}

\begin{tcolorbox}[colback=lightgray,colframe=secondary,title=\textbf{Experimento 2.1: Probar el modelo con datos pasados (Backtesting)}]
\begin{itemize}[leftmargin=*]
    \item \textbf{Lo que queremos saber:} ¿Nuestro modelo de Inteligencia Artificial realmente predice mejor el pago de independientes que los modelos de scoring tradicionales?
    \item \textbf{Cómo lo probamos:} Tomamos una base de datos anónima de préstamos pasados y comparamos los resultados del score tradicional frente a las predicciones de nuestro algoritmo.
    \item \textbf{Por qué importa la métrica:} Demostrar una reducción del 25\% en los rechazos injustificados (\textit{falsos negativos}) significa que la entidad financiera puede aprobar a miles de clientes independientes adicionales que sí pueden pagar, aumentando sus ingresos sin incrementar la mora.
\end{itemize}
\end{tcolorbox}

\begin{tcolorbox}[colback=lightgray,colframe=secondary,title=\textbf{Experimento 2.2: Prueba de usabilidad de reportes con analistas de riesgo}]
\begin{itemize}[leftmargin=*]
    \item \textbf{Lo que queremos saber:} ¿Los reportes explicativos visuales ayudan a los analistas a tomar decisiones de crédito más rápido?
    \item \textbf{Cómo lo probamos:} Le mostramos dos tipos de pantalla a 10 analistas de crédito (una con solo un puntaje numérico y otra con el desglose explicativo de StatCredit AI) y tomamos el tiempo de evaluación de cada caso.
    \item \textbf{Por qué importa la métrica:} Pasar de 15 minutos a menos de 7 minutos por cliente reduce a la mitad el tiempo de atención. Esto permite al equipo de analistas procesar el doble de solicitudes al día sin necesidad de contratar más personal.
\end{itemize}
\end{tcolorbox}

\subsubsection{Pregunta 3: ¿Cómo adquiere tu cliente tu producto?}

\begin{tcolorbox}[colback=lightgray,colframe=secondary,title=\textbf{Experimento 3.1: Prueba de facilidad de integración de la API}]
\begin{itemize}[leftmargin=*]
    \item \textbf{Lo que queremos saber:} ¿Es fácil y rápido para los equipos de sistemas conectar nuestra API de scoring a sus plataformas?
    \item \textbf{Cómo lo probamos:} Invitamos a 5 programadores de fintechs aliadas a conectar nuestra API en un ambiente de pruebas (\textit{Sandbox}) y medimos el tiempo que tardan en realizar la primera consulta exitosa.
    \item \textbf{Por qué importa la métrica:} Si logran conectar el servicio en menos de 1 hora, eliminamos la fricción técnica y la resistencia del área de tecnología, que suele ser uno de los mayores obstáculos para cerrar ventas B2B.
\end{itemize}
\end{tcolorbox}

\begin{tcolorbox}[colback=lightgray,colframe=secondary,title=\textbf{Experimento 3.2: Prueba de versión gratuita inicial (Freemium Pilot)}]
\begin{itemize}[leftmargin=*]
    \item \textbf{Lo que queremos saber:} ¿Regalar un paquete de 500 consultas iniciales ayuda a que las empresas decidan contratar el servicio pagado?
    \item \textbf{Cómo lo probamos:} Ofrecemos un plan piloto sin costo de 500 consultas y medimos qué porcentaje de empresas deciden firmar un contrato mensual cuando se agota la prueba.
    \item \textbf{Por qué importa la métrica:} 500 consultas le permiten a la fintech probar la herramienta con clientes reales suficientes. Lograr que al menos 2 de cada 10 empresas (20\%) se conviertan a clientes de pago confirma que la prueba gratuita genera suficiente confianza comercial.
\end{itemize}
\end{tcolorbox}

\subsubsection{Pregunta 4: ¿Cómo ganas dinero con tu producto?}

\begin{tcolorbox}[colback=lightgray,colframe=secondary,title=\textbf{Experimento 4.1: Encuesta de preferencia de modelo de cobro}]
\begin{itemize}[leftmargin=*]
    \item \textbf{Lo que queremos saber:} ¿Los clientes prefieren pagar una mensualidad fija por el software o un cobro por cada consulta de score realizada?
    \item \textbf{Cómo lo probamos:} Realizamos encuestas sencillas de decisión a encargados de compra en entidades crediticias, ofreciéndoles distintas combinaciones de tarifas.
    \item \textbf{Por qué importa la métrica:} Descubrir que prefieren un cobro híbrido (una cuota mensual base baja más un valor pequeño por consulta realizada) nos ayuda a fijar un precio competitivo que cubre nuestros costos fijos y crece según el volumen del cliente.
\end{itemize}
\end{tcolorbox}

\begin{tcolorbox}[colback=lightgray,colframe=secondary,title=\textbf{Experimento 4.2: Simulación de estabilidad financiera (Monte Carlo)}]
\begin{itemize}[leftmargin=*]
    \item \textbf{Lo que queremos saber:} ¿El modelo de negocio de StatCredit AI sigue siendo rentable si algunos clientes cancelan o si los costos de los servidores aumentan?
    \item \textbf{Cómo lo probamos:} Realizamos una simulación matemática (Monte Carlo) cambiando variables de ingresos, costos de servidores e inflación en 10,000 escenarios posibles.
    \item \textbf{Por qué importa la métrica:} Si el negocio mantiene ganancias en más del 90\% de los escenarios simulados, comprobamos que el proyecto es financieramente estable y resistente a imprevistos del mercado.
\end{itemize}
\end{tcolorbox}

\subsubsection{Pregunta 5: ¿Cómo diseñas y construyes tu producto?}

\begin{tcolorbox}[colback=lightgray,colframe=secondary,title=\textbf{Experimento 5.1: Prueba de velocidad y carga del sistema}]
\begin{itemize}[leftmargin=*]
    \item \textbf{Lo que queremos saber:} ¿La API responde a tiempo cuando entran cientos de solicitudes al mismo segundo?
    \item \textbf{Cómo lo probamos:} Enviamos 500 solicitudes por segundo simuladas a nuestros servidores para medir el tiempo de respuesta y verificar que no se produzcan caídas del servicio.
    \item \textbf{Por qué importa la métrica:} En aplicaciones móviles de crédito inmediato, si la respuesta tarda más de 0.3 segundos (300 milisegundos), los usuarios abandonan la solicitud en el punto de pago. Garantizar respuestas en menos de 0.3 segundos es vital para la experiencia del usuario.
\end{itemize}
\end{tcolorbox}

\begin{tcolorbox}[colback=lightgray,colframe=secondary,title=\textbf{Experimento 5.2: Prueba piloto a ciegas con entidad financiera}]
\begin{itemize}[leftmargin=*]
    \item \textbf{Lo que queremos saber:} ¿Las personas independientes aprobadas con nuestro score realmente pagan a tiempo en la vida real?
    \item \textbf{Cómo lo probamos:} Evaluamos 500 solicitudes reales en paralelo con una entidad aliada a ciegas (el banco evalúa con su método tradicional y nosotros con StatCredit AI, sin interferir en la decisión final).
    \item \textbf{Por qué importa la métrica:} Una tasa de mora inferior al 3\% en el grupo evaluado por nuestro sistema después de 90 días demuestra que la Inteligencia Artificial funciona en el mercado real y cumple con los parámetros de riesgo permitidos por la regulación (SARC).
\end{itemize}
\end{tcolorbox}

\subsubsection{Pregunta 6: ¿Cómo escalas tu negocio?}

\begin{tcolorbox}[colback=lightgray,colframe=secondary,title=\textbf{Experimento 6.1: Prueba de adaptación del algoritmo a otro país (México)}]
\begin{itemize}[leftmargin=*]
    \item \textbf{Lo que queremos saber:} ¿El modelo entrenado con datos de independientes en Colombia se puede usar en otro país como México realizando ajustes mínimos?
    \item \textbf{Cómo lo probamos:} Evaluamos nuestro algoritmo utilizando un conjunto de datos transaccionales de micronegocios e independientes en México.
    \item \textbf{Por qué importa la métrica:} Mantener más del 80\% de la precisión inicial demuestra que la lógica de comportamiento financiero informal es similar entre países latinoamericanos, lo que permite expandir el negocio a nuevos mercados rápidamente y sin grandes costos de desarrollo.
\end{itemize}
\end{tcolorbox}

\begin{tcolorbox}[colback=lightgray,colframe=secondary,title=\textbf{Experimento 6.2: Inclusión de datos de facturación electrónica (DIAN)}]
\begin{itemize}[leftmargin=*]
    \item \textbf{Lo que queremos saber:} ¿Agregar datos de facturas electrónicas mejora la capacidad de predecir el pago en micronegocios?
    \item \textbf{Cómo lo probamos:} Comparamos la precisión del modelo evaluando primero solo datos bancarios y luego agregando datos de ventas y compras en facturas electrónicas.
    \item \textbf{Por qué importa la métrica:} Las facturas reflejan las ventas reales recientes de un pequeño negocio. Un aumento del 15\% en la precisión justifica crear un módulo especializado para micronegocios, ampliando nuestro catálogo de productos.
\end{itemize}
\end{tcolorbox}

\subsection{Conclusión de la Tarea 2}
Estos 12 experimentos le brindan a \textbf{StatCredit AI} una estrategia clara de prueba y aprendizaje. Nos permiten validar paso a paso nuestro producto con datos reales, asegurando que cada funcionalidad que desarrollemos responda a una necesidad comprobada del mercado financiero en Colombia.

\newpage
% ---------------------------------------------------------
% SECCIÓN 3: TAREA 3
% ---------------------------------------------------------
\section{Tarea 3: Business Model Canvas}

A continuación se presenta el Business Model Canvas estructurado para \textbf{StatCredit AI}, desglosando los 9 componentes del modelo de negocio:

\begin{table}[h!]
\centering
\small
\begin{tabularx}{\linewidth}{|X|X|X|X|X|}
\hline
\rowcolor{lightgray}
\textbf{8. Alianzas Clave} & 
\textbf{7. Actividades Clave} & 
\textbf{2. Propuesta de Valor} & 
\textbf{4. Relaciones con Clientes} & 
\textbf{1. Segmentos de Clientes} \\ 
\hline
\begin{itemize}[leftmargin=*]
    \item Plataformas de Datos (Belvo).
    \item Billeteras digitales (Nequi, Daviplata).
    \item Gremios (Colombia Fintech, Fecolfin).
    \item Servidores cloud (AWS / GCP).
\end{itemize} & 
\begin{itemize}[leftmargin=*]
    \item Entrenar modelos de IA causales.
    \item Mantener la API rápida y 24/7.
    \item Auditar reportes ante norma SARC.
\end{itemize} & 
\begin{itemize}[leftmargin=*]
    \item \textbf{Score Alternativo:} Evalúa capacidad de repago real para aumentar aprobaciones.
    \item \textbf{Reportes Explicables:} Muestra por qué la IA toma cada decisión.
    \item \textbf{API Rápida:} Conexión sencilla e integración en pocas horas.
\end{itemize} & 
\begin{itemize}[leftmargin=*]
    \item Acompañamiento técnico directo.
    \item Garantía de alta disponibilidad del servicio.
    \item Asesoría en comités y auditorías.
\end{itemize} & 
\begin{itemize}[leftmargin=*]
    \item \textbf{Primer Cliente (Beachhead):} Fintechs y neobancos.
    \item \textbf{Segundo Grupo:} Cooperativas de ahorro y crédito.
    \item \textbf{Tercer Grupo:} Banca tradicional.
\end{itemize} \\
\cline{2-2} \cline{4-4}
 & \rowcolor{lightgray} 
\textbf{6. Recursos Clave} & & \rowcolor{lightgray} 
\textbf{3. Canales} & \\
\cline{2-2} \cline{4-4}
 & \begin{itemize}[leftmargin=*]
    \item Algoritmos de IA explicable (SHAP).
    \item Equipo de Ciencia de Datos.
    \item Datos históricos de prueba.
\end{itemize} & & \begin{itemize}[leftmargin=*]
    \item Venta directa B2B a riesgos.
    \item Portal web con prueba gratis API.
    \item AWS Marketplace y ferias fintech.
\end{itemize} & \\
\hline
\end{tabularx}
\end{table}

\begin{table}[h!]
\centering
\small
\begin{tabularx}{\linewidth}{|X|X|}
\hline
\rowcolor{lightgray}
\textbf{9. Estructura de Costos} & 
\textbf{5. Fuentes de Ingresos} \\
\hline
\begin{itemize}[leftmargin=*]
    \item Servidores cloud e infraestructura GPU (AWS/GCP).
    \item Sueldos del equipo de Ciencia de Datos y desarrollo.
    \item Costos de consumo de APIs de agregación Open Finance.
\end{itemize} & 
\begin{itemize}[leftmargin=*]
    \item Suscripción mensual base (cuota fija) + costo variable por consulta de score exitosa.
    \item Licencias anuales Enterprise para bancos grandes.
    \item Módulos de reportes explicativos avanzados.
\end{itemize} \\
\hline
\end{tabularx}
\end{table}

\subsection{Conclusión del Portafolio}
El consolidado del \textbf{Portafolio de Evidencias No. 2} demuestra que \textbf{StatCredit AI} cuenta con una ruta clara de validación, experimentos concretos y un modelo de negocio sostenible para transformar el acceso al crédito en Colombia.

\end{document}
